\documentclass[10pt,letterpaper]{article}
\usepackage[T1]{fontenc}
\usepackage{courier}
\renewcommand{\familydefault}{\ttdefault}
\usepackage{amsmath,amssymb,amsthm}
\usepackage[margin=0.78in]{geometry}
\usepackage{microtype}
\usepackage{fancyhdr}
\usepackage{titlesec}
\usepackage{xcolor}
\usepackage{hyperref}

\definecolor{ink}{HTML}{232323}
\hypersetup{colorlinks=true,linkcolor=ink,urlcolor=ink,
  pdftitle={Can shared references sharpen region recovery?},
  pdfauthor={Jonathan R. Landers}}
\color{ink}
\setlength{\parindent}{0pt}
\setlength{\parskip}{0.50em}
\setlength{\headheight}{14pt}
\linespread{1.03}
\emergencystretch=5em
\pagestyle{fancy}
\fancyhf{}
\fancyhead[L]{\small Distribution Sampling Neighborhoods}
\fancyhead[R]{\small 18 OCT 2026}
\fancyfoot[R]{\small \thepage}
\renewcommand{\headrulewidth}{0.35pt}
\renewcommand{\footrulewidth}{0pt}
\newcommand{\sectionrule}{%
  \par\nobreak\vspace{0.12em}\noindent
  \makebox[\textwidth][c]{\textcolor{black}{\rule{0.75\textwidth}{0.35pt}}}}
\titleformat{\section}{\normalfont\bfseries\large}{\thesection.}{0.6em}{}[\sectionrule]
\titlespacing*{\section}{0pt}{0.85em}{0.30em}
\newtheoremstyle{typed}{0.5em}{0.5em}{\normalfont}{}{\bfseries}{.}{0.5em}{}
\theoremstyle{typed}
\newtheorem{lemma}{Lemma}
\newtheorem{proposition}[lemma]{Proposition}
\renewcommand{\proofname}{\textbf{\upshape Proof}}

\begin{document}
\thispagestyle{fancy}

{\Large\bfseries Can shared references sharpen region recovery?}\par
{\small A finite-sample extension of the inverse sampling model}\par
{\small Jonathan R. Landers}\par
\vspace{0.25em}\noindent\rule{\textwidth}{0.5pt}

\section{From relative shape to a shared space}

The inverse sampling note~\cite{inverse} shows that pairwise distances
within a constrained sample can leave its location and original mass
ambiguous. This note adds shared reference points and a restricted family
of possible regions, then bounds the mass recovered incorrectly.

At a fixed input, let $P_0$ be the original law on $\mathbb R^d$ and let
$B$ be a measurable region with $\beta=P_0(B)>0$. As before, assume hard
conditioning and $n\geq1$ independent constrained outputs:
\begin{equation}
  Q(A)=\frac{P_0(A\cap B)}{\beta},\qquad Y_1,\ldots,Y_n\sim Q.
  \label{eq:model}
\end{equation}
Original samples can serve as anchors $a_0,\ldots,a_d$. Measure the
distance from every output to the \emph{same} anchors, retaining their
identities. This requires distances between the two datasets. The
shared representation is
\begin{equation}
  \phi(y)=\bigl(\|y-a_0\|,\ldots,\|y-a_d\|\bigr).
  \label{eq:anchors}
\end{equation}

\begin{lemma}[Exact location from anchors]
If the anchors are affinely independent, $\phi$ is one-to-one.
Known anchor coordinates and exact distances recover $y$ by a linear solve.
\end{lemma}
\begin{proof}
Write $r_j=\|y-a_j\|$. Subtracting squared distances gives, for $j=1,\ldots,d$,
\[
  2(a_j-a_0)^\top y
  =\|a_j\|^2-\|a_0\|^2+r_0^2-r_j^2.
\]
The $d$ row vectors $(a_j-a_0)^\top$ are linearly independent.
The system therefore has a unique solution. Identical distance vectors
must correspond to the same point.
\end{proof}

Thus three noncollinear anchors suffice in the plane. The construction
is a distance-reconstruction method~\cite{edm}; an arbitrary metric or
degenerate anchor set need not have this property. It locates observed
points exactly, while unsampled parts of $B$ still require inference.

\section{A small family of possible regions}

Fix a finite family $\mathcal H=\{C_1,\ldots,C_M\}$ of measurable regions,
with $B\in\mathcal H$. Candidates may be defined in the shared space and
pulled back through $\phi$. Fix the anchors and family before collecting
the samples used below; an independent pilot sample is allowed.

\newpage
\section{Minimum mass gives an error bound}

When candidate masses are known, retain candidates containing all
constrained outputs and choose the one with least original mass:
\begin{equation}
  \widehat B\in\operatorname*{arg\,min}_{C\in\mathcal H:\,Y_i\in C\ \forall i}
  P_0(C).
  \label{eq:estimator}
\end{equation}
A fixed ordering resolves ties. The true $B$ is feasible almost surely,
so the search is nonempty. Under~\eqref{eq:model}, the likelihood relative
to $P_0^{\otimes n}$ is
\[
  L(C)=P_0(C)^{-n}\prod_{i=1}^n\mathbf1\{Y_i\in C\}.
\]
Equation~\eqref{eq:estimator} maximizes this likelihood. Zero-mass
candidates are infeasible almost surely and can be discarded. This
extends the original note's two-interval rule and is related to
minimum-volume set estimation~\cite{scott}.

\begin{proposition}[Known original masses]
For $0<\delta<1$, with probability at least $1-\delta$,
\begin{equation}
  \boxed{\quad P_0(B\triangle\widehat B)
    \leq \frac{2\beta}{n}\log\frac{M}{\delta}.\quad}
  \label{eq:exactbound}
\end{equation}
Here $B\triangle C=(B\setminus C)\cup(C\setminus B)$ counts both missed
and incorrectly included original mass; $\log$ is the natural logarithm.
\end{proposition}
\begin{proof}
For a fixed candidate $C$, put
$q_C=Q(B\setminus C)=P_0(B\setminus C)/\beta$. Independence gives
\[
  \Pr\{Y_i\in C\ \forall i\}=(1-q_C)^n\leq e^{-nq_C}.
\]
Taking a union bound over the $M$ candidates, for every $t>0$,
\[
 \Pr\{\text{some feasible }C\text{ has }q_C>t\}\leq M e^{-nt}.
\]
Set $t=\log(M/\delta)/n$. Except on an event of probability at most
$\delta$, every feasible candidate, including $\widehat B$, misses
at most $\beta t$ original mass.

Write $a=P_0(B\setminus\widehat B)$ and
$b=P_0(\widehat B\setminus B)$. Minimality and feasibility of $B$ give
\[
 P_0(\widehat B)=\beta-a+b\leq\beta,
 \qquad\text{hence }b\leq a.
\]
Therefore $P_0(B\triangle\widehat B)=a+b\leq2a\leq2\beta t$,
which proves~\eqref{eq:exactbound}.
\end{proof}

The two steps have distinct roles: observing every constrained output
limits missed mass; minimizing original mass also limits extra mass.
For a desired error $\varepsilon>0$, it suffices that
$n\geq(2\beta/\varepsilon)\log(M/\delta)$. Replacing $\beta$ by $1$
gives a conservative guarantee when the true mass is unknown.

\newpage
\section{When the original law is available by sampling}

Draw $m\geq1$ independent outputs $X_1,\ldots,X_m\sim P_0$, separately
from the constrained and pilot samples. Replace $P_0(C)$ in~\eqref{eq:estimator} by
$\widehat P_m(C)=m^{-1}\sum_{j=1}^m\mathbf1\{X_j\in C\}$, and call the
result $\widetilde B$. Original samples supply mass estimates; they do
not supply membership labels for $B$.

\begin{proposition}[Estimated original masses]
Under the same fixed-family assumption, with probability at least
$1-\delta$,
\begin{equation}
 \boxed{\quad P_0(B\triangle\widetilde B)
 \leq \frac{2\beta}{n}\log\frac{2M}{\delta}
       +\sqrt{\frac{2}{m}\log\frac{4M}{\delta}}.\quad}
 \label{eq:samplebound}
\end{equation}
\end{proposition}
\begin{proof}
Hoeffding's inequality for each indicator average, followed by a union
bound, gives
\[
 \Pr\!\left\{\max_{C\in\mathcal H}|\widehat P_m(C)-P_0(C)|>\eta\right\}
 \leq 2M e^{-2m\eta^2}.
\]
Choose $\eta=\sqrt{\log(4M/\delta)/(2m)}$, making this probability at
most $\delta/2$. On the complementary event, empirical minimality gives
\[
 P_0(\widetilde B)\leq\widehat P_m(\widetilde B)+\eta
 \leq\widehat P_m(B)+\eta\leq\beta+2\eta.
\]
The preceding proof, with failure probability $\delta/2$, also gives
$a=P_0(B\setminus\widetilde B)\leq\beta\log(2M/\delta)/n$ for every
feasible choice, including $\widetilde B$. Both events hold with
probability at least $1-\delta$. Writing
$b=P_0(\widetilde B\setminus B)$, the mass comparison implies
$b\leq a+2\eta$. Thus the total error is at most $2a+2\eta$,
which is~\eqref{eq:samplebound}.
\end{proof}

\section{What the extension buys}

More original samples shrink the second term; more constrained samples
shrink the first. A shared representation can help if it makes a small,
correct candidate family available. An invertible change of coordinates
alone preserves probability error. Guessing an alignment does not
resolve indistinguishable alternatives without additional information.

These are finite-family guarantees, not a claim of tighter numerical
error than the original experiment. Arbitrary two-interval endpoints
form an infinite family; a finite grid may exclude the true region.
Misspecification, noisy anchors, or a family learned from the inference
samples require additional analysis. Exact sample locations do not
certify exact unseen boundaries.

{\footnotesize\raggedright
\begin{thebibliography}{3}
\setlength{\itemsep}{0pt}
\bibitem{inverse} J. R. Landers.
\href{inverse_sampling_region.pdf}{Can sampled distances recover a constrained region?}
Companion note, 2026.
\bibitem{edm} I. Dokmani\'{c}, R. Parhizkar, J. Ranieri, and M. Vetterli.
\href{https://arxiv.org/abs/1502.07541}{Euclidean distance matrices: Essential theory, algorithms and applications.}
2015.
\bibitem{scott} C. D. Scott and R. D. Nowak.
\href{https://jmlr.org/papers/v7/scott06a.html}{Learning minimum volume sets.}
JMLR 7:665--704, 2006.
\end{thebibliography}}
\end{document}
